% preamble.tex
\usepackage[T2A]{fontenc}
\usepackage[utf8]{inputenc}
\usepackage[russian]{babel}
\usepackage{amsmath}
\usepackage{amssymb}
\usepackage{graphicx}
\usepackage{float}
\usepackage{geometry}

% === Для графиков (pgfplots/TikZ) ===
\usepackage{pgfplots}
\usetikzlibrary{arrows.meta, decorations.pathmorphing, intersections}
\pgfplotsset{compat=1.18}
% ====================================

\geometry{a4paper, margin=1in}

% === Теоремные окружения с общей нумерацией в пределах section ===
\newtheorem{theorem}{Теорема}[section]
\newtheorem{lemma}[theorem]{Лемма}
\newtheorem{property}[theorem]{Свойство}
\newtheorem{definition}[theorem]{Определение}
\newtheorem{corollary}[theorem]{Следствие}
\newtheorem{remark}[theorem]{Замечание}
\newtheorem{example}[theorem]{Пример}
\newtheorem{proof}[theorem]{Доказательство}
% =================================================

% Настройка гиперссылок
\usepackage{hyperref}
\hypersetup{
    colorlinks=true,
    linkcolor=blue,
    citecolor=blue,
    urlcolor=blue,
    pdftitle={Основной экз 2сем линал Камынин}
}

